Valuing a shock and recovering its effects from data are different tasks. Initial
expenditure shares value price changes once those changes and income responses
are known. Aggregate welfare results in Arkolakis, Costinot and Rodriguez-Clare
(2012) instead rely on restrictions defining the class of trade models they
consider. Their result is useful for comparing what information is sufficient
within that class, rather than as a general proof of household valuation under
unrestricted preferences. The hypothetical tourist application starts from local
valuation and then asks how to recover its income and price terms.

That recovery question brings demand restrictions into the argument. Redding and
Weinstein (2020) study price-index recovery under CES preferences, including how
demand residuals and normalization affect the recovered index. A destination
effect in a spending equation is therefore not automatically a measured price
change. Its interpretation depends on the demand system and the chosen
normalization. This connection helps organize the correction from spending to
income, but it does not establish that the tourist-demand variation is exogenous.
The latter requires a separate condition on other influences on resident spending.

Spatial interactions raise a further question: which responses occur away from
the directly affected destination? Allen, Arkolakis and Takahashi (2020) provide
a gravity representation, while Allen, Arkolakis and Li (2024) characterize
equilibrium properties under stated conditions on interactions. Adao, Arkolakis
and Esposito, in their September 2025 manuscript, use a spatial response
representation to distinguish differential and aggregate effects of trade shocks.
These connections concern different parts of the same problem: representing
links, characterizing the resulting system, and relating a response to an empirical
specification. The tourist application uses work and purchase links to recover a
local welfare response under its maintained restrictions. It does not inherit
causal validity from an equilibrium theorem or solve a general dynamic economy.
Keeping these tasks separate makes the literature comparison useful for deciding
what the observed purchases can tell us.
